\section{安全保障貿易管理：仕様・相手・用途を確認する}
\label{theory:export:screening}
同じ測定器が民生研究にも軍事目的にも役立つ場合、買い手の肩書だけで用途は分からない。安全保障貿易管理は、軍事転用可能な貨物や技術が懸念活動へ渡ることを防ぐため、外為法等に基づいて行う\parencite{export_overview}。

まず貨物や技術の仕様を法令のリストへ照合する該非判定を行う。次に、相手・需要者・用途・仕向地等から取引を審査する。リスト非該当でも、補完的なキャッチオール規制の条件に該当する場合がある。最後に、審査した物と発送する物の同一性、許可の取得と内容を出荷時に確認する。少量サンプル、国内で買える製品、自作装置、手荷物での持出しというだけで、この確認を省けない。

メーカーの判定書は重要な資料だが、古い判定書が現在の型式・仕様・法令に合うかを輸出者も確かめる。信頼する取引先からの紹介も、新たな需要者の用途確認を代替しない。営業・技術・管理・出荷の各部門で判断根拠を共有することで、適切な審査を実際の発送へ結び付ける。

\section{技術提供と例外：場所より情報の内容を調べる}
\label{theory:export:exceptions}
技術には設計・製造・使用に必要な情報を、データや助言として提供する場合がある。国内の非居住者への提供も対象となり得るので、日本国内であることだけで規制外と判断しない\parencite{export_scope2022}。大学の非公開な共同研究打合せで技術を交換することと、不特定多数が入手できる公知技術の発表は違う。

装置の部分品特例も、部品が組み込まれているという一事だけで使えるわけではない。公的事例集は、ロボットに内蔵されたセンサーが主要要素でないことを学内の担当部署が確認する手順を示す\parencite{export_cases}。特例の対象と仕様・価格等の条件を検討し、本人だけの思い込みで適用しない。特許出願等に必要な最小限の技術提供についての例外も、手続目的を越えるノウハウ提供まで含めない。当時条文の直接照合は初稿の留保事項である。

\begin{tcolorbox}[title=出題時点と現在の規則]
ここでの過去問は2021--2023年の事例である。経済産業省の現在の案内には2025年10月9日施行の補完的規制見直しが含まれる\parencite{export_current_updates}。現在の品目・対象国・用途・通知条件を過去問へ補って判定しない。取引の実施では、実施時点の規則を確認する。
\end{tcolorbox}
